% theory_fundamentals.tex —— 潮流计算 通用理论：节点方程、Newton 法与数值性态
% 本文件为理论层章节，遵循 docs/modules/THEORY_WRITING_GUIDE.md。

\section{潮流计算的数学基础}\label{sec:pf-th-fund}

\begin{theorynote}
本章为通用数学理论，按电力系统潮流计算的标准模型与数值分析文献体系撰写，覆盖：
节点功率方程从基尔霍夫电流定律与 $\pi$ 等值电路出发的一般推导（含复数变比变压器
与移相器两端口）、潮流问题的数学分类（PQ/PV/V$\theta$ 节点、方程-未知数计数、解的
存在性与多解性、可解性的压缩充分条件）、Newton--Raphson 法的一般理论（局部二次收敛定理及其证明、收敛域、
雅可比奇异的物理含义）、直角坐标下失配的精确二次结构与 Iwamoto 最优乘子法、
病态与数值性态（条件数分析、高 $R/X$ 比、重负荷/弱网、
雅可比奇异与电压崩溃的关系）、稀疏技术理论（图消元与填充、最小度/近似最小度
排序原理、LU 填充的图刻画），以及由平台求解器实测数据给出的收敛性数值算例与基准。
与平台实现（\srcpath{src/power\_flow/}）的对应关系见
本章末"与实现的对应"表，超出实现的内容以"未实现扩展说明"框就地标记。
\end{theorynote}

\subsection{记号与约定}
\begin{itemize}
  \item $\mathcal N=\{1,\dots,n\}$：AC 节点集；$n$ 为节点数；$i,j$ 为节点下标；
  \item $V_i=V_{m,i}e^{\mathrm j\theta_i}\in\mathbb C$：节点 $i$ 的复电压，
        $\theta_{ij}=\theta_i-\theta_j$；$\overline{(\cdot)}$ 表示复共轭；
        $\mathrm j$ 为虚数单位；
  \item $S_i=P_i+\mathrm jQ_i=V_i\,\overline{I_i}$：节点注入复功率（注入网络为正方向）；
  \item $Y_{bus}=G+\mathrm jB\in\mathbb C^{n\times n}$：节点导纳矩阵，
        $Y_{ij}=G_{ij}+\mathrm jB_{ij}$；
  \item 支路 $i$--$j$：串联阻抗 $z_{ij}=r_{ij}+\mathrm jx_{ij}$，串联导纳
        $y_{ij}=1/z_{ij}=g_{ij}+\mathrm jb_{ij}$，线路充电（对地）总电纳 $b_{c,ij}$，
        两端各半；
  \item $t=ae^{\mathrm j\varphi}$：变压器/移相器复数变比，$a>0$ 为变比幅值，
        $\varphi$ 为移相角；约定理想变压器位于 from 侧（节点 $i$ 侧）；
  \item $F:\mathbb R^m\to\mathbb R^m$：潮流失配（mismatch）函数，$J=DF$ 为其雅可比；
        $\lVert\cdot\rVert$ 未指明时为向量 $2$-范数及其诱导矩阵范数；
        $\kappa(J)=\lVert J\rVert\,\lVert J^{-1}\rVert$：条件数；
  \item $u$：浮点运算单位舍入（双精度 $u\approx1.1\times10^{-16}$）；
  \item 全部电气量除特别声明外取标幺值（系统基准 $S_B$，电压基准按节点额定值）。
\end{itemize}

\subsection{节点功率方程的一般推导}\label{sec:pf-th-fund-power-eq}

\subsubsection{\texorpdfstring{$\pi$}{π} 等值支路与复变比两端口}

\begin{definition}[$\pi$ 等值支路]\label{def:pf-th-fund-pi}
输电线路/电缆支路 $i$--$j$ 的集总参数 $\pi$ 等值电路由串联导纳 $y_{ij}$ 与两端
对地充电电纳 $\mathrm jb_{c,ij}/2$ 组成。记 $i$、$j$ 端注入支路的电流为
$I_{ij}$、$I_{ji}$（流入支路为正），则无变压器支路满足
\begin{equation}
\begin{bmatrix}I_{ij}\\ I_{ji}\end{bmatrix}
=
\begin{bmatrix}
y_{ij}+\mathrm jb_{c,ij}/2 & -y_{ij}\\
-y_{ij} & y_{ij}+\mathrm jb_{c,ij}/2
\end{bmatrix}
\begin{bmatrix}V_i\\ V_j\end{bmatrix}.
\label{eq:pf-th-fund-pi}
\end{equation}
该两端口是互易的（转移导纳相等）且无源，$2\times2$ 导纳矩阵对称。
\end{definition}

\begin{definition}[复变比变压器/移相器两端口]\label{def:pf-th-fund-tap}
非标准变比变压器与移相器统一建模为：$\pi$ 等值串联支路级联一台位于 from 侧
（节点 $i$ 侧）的理想复变比变压器。设理想变压器内侧（绕组侧）节点为 $k$，复数变比
$t=ae^{\mathrm j\varphi}$ 定义为
\begin{equation}
V_k=\frac{V_i}{t}.
\end{equation}
理想变压器不储能、不耗能，两侧复功率守恒：
\begin{equation}
V_i\,\overline{I_i}=V_k\,\overline{I_k'}
\quad\Longrightarrow\quad
I_i=\frac{I_k'}{\overline{t}},
\label{eq:pf-th-fund-ideal}
\end{equation}
其中 $I_k'$ 为内侧节点 $k$ 流入串联支路的电流。
\end{definition}

\begin{proposition}[含变比与移相的一般支路模型]\label{prop:pf-th-fund-tport}
定义~\ref{def:pf-th-fund-tap} 的两端口导纳参数为
\begin{equation}
\begin{bmatrix}I_i\\ I_j\end{bmatrix}
=
\begin{bmatrix}
\dfrac{y_{ij}+\mathrm jb_{c,ij}/2}{|t|^2} & -\dfrac{y_{ij}}{\overline{t}}\\[2.2ex]
-\dfrac{y_{ij}}{t} & y_{ij}+\mathrm jb_{c,ij}/2
\end{bmatrix}
\begin{bmatrix}V_i\\ V_j\end{bmatrix}.
\label{eq:pf-th-fund-tport}
\end{equation}
\end{proposition}

\begin{proof}
内侧节点 $k$ 经 $\pi$ 支路与节点 $j$ 相连，由式~\eqref{eq:pf-th-fund-pi}，
\begin{equation}
I_k'=\Bigl(y_{ij}+\mathrm j\tfrac{b_{c,ij}}{2}\Bigr)V_k-y_{ij}V_j,\qquad
I_j=-y_{ij}V_k+\Bigl(y_{ij}+\mathrm j\tfrac{b_{c,ij}}{2}\Bigr)V_j .
\end{equation}
代入 $V_k=V_i/t$ 与式~\eqref{eq:pf-th-fund-ideal} 的 $I_i=I_k'/\overline{t}$：
\begin{equation}
I_i=\frac{1}{\overline{t}}\Bigl[\Bigl(y_{ij}+\mathrm j\tfrac{b_{c,ij}}{2}\Bigr)
\frac{V_i}{t}-y_{ij}V_j\Bigr]
=\frac{y_{ij}+\mathrm jb_{c,ij}/2}{|t|^2}V_i-\frac{y_{ij}}{\overline{t}}V_j ,
\end{equation}
\begin{equation}
I_j=-\frac{y_{ij}}{t}V_i+\Bigl(y_{ij}+\mathrm j\tfrac{b_{c,ij}}{2}\Bigr)V_j .
\end{equation}
两式合并即得式~\eqref{eq:pf-th-fund-tport}。
\end{proof}

\begin{remark}[互易性、对称性与移相角的作用]\label{rem:pf-th-fund-reciprocity}
当 $\varphi=0$（实变比）时 $-\,y_{ij}/\overline{t}=-\,y_{ij}/t$，两端口互易，
$Y_{bus}$ 保持（复）对称；当 $\varphi\neq0$ 时
$Y_{ij}^{ft}\neq Y_{ij}^{tf}$，两端口在导纳参数意义下不再互易，$Y_{bus}$
失去对称性——但理想变压器仍满足功率守恒式~\eqref{eq:pf-th-fund-ideal}，
非对称恰恰是"电压相移而功率守恒"在导纳框架中的正确编码。工程上
$\varphi\neq0$ 的支路（移相器、相位角调节变压器）用于控制有功潮流的分布：
其效果是在支路两端电压相角差上叠加一个可控偏置 $\varphi$。任何假设
$Y_{bus}$ 对称的求解器或求逆算法在含移相器的网络上必须先核对适用性。
\end{remark}

\subsubsection{节点导纳矩阵与节点功率方程}

\begin{definition}[节点导纳矩阵]\label{def:pf-th-fund-ybus}
对所有在运支路按式~\eqref{eq:pf-th-fund-pi}/\eqref{eq:pf-th-fund-tport}
组装：对角元累加自导纳（含充电与节点对地并联导纳 $g_{s,i}+\mathrm jb_{s,i}$），
非对角元累加转移导纳，得 $Y_{bus}=G+\mathrm jB$。其结构性质：
(i) 稀疏——非零元个数约为 $n+2b$（$b$ 为支路数，电力网 $b\approx1.5n$）；
(ii) 无移相器时对称；（iii) 一般\textbf{不}具对角占优性（充电电纳与串联电纳
在对角元中异号抵消），因此不能直接套用对角占优类的迭代法理论。
\end{definition}

\begin{theorem}[节点功率方程]\label{thm:pf-th-fund-power-eq}
设网络无内部独立电流源，节点注入电流仅由支路与对地导纳产生，则
\begin{equation}
I_i=\sum_{j\in\mathcal N}Y_{ij}V_j,\qquad
S_i=V_i\,\overline{\Bigl(\sum_{j\in\mathcal N}Y_{ij}V_j\Bigr)},
\qquad i\in\mathcal N .
\label{eq:pf-th-fund-si}
\end{equation}
\end{theorem}

\begin{proof}
对节点 $i$ 应用基尔霍夫电流定律：注入网络的电流 $I_i$ 等于经各关联支路流出
电流之和。每条关联支路按式~\eqref{eq:pf-th-fund-pi}/\eqref{eq:pf-th-fund-tport}
的右端贡献一个关于端电压的线性式；把所有支路的端电压系数按列下标归类，
恰为 $Y_{bus}$ 第 $i$ 行与电压向量的乘积，即第一式。第二式由复功率定义
$S_i=V_i\overline{I_i}$ 直接代入。
\end{proof}

\begin{corollary}[极坐标功率方程]\label{cor:pf-th-fund-polar}
写 $V_i=V_{m,i}e^{\mathrm j\theta_i}$、$Y_{ij}=G_{ij}+\mathrm jB_{ij}$，
$\theta_{ij}=\theta_i-\theta_j$，则式~\eqref{eq:pf-th-fund-si} 的实部、虚部为
\begin{align}
P_i&=\sum_{j\in\mathcal N}V_{m,i}V_{m,j}
\bigl(G_{ij}\cos\theta_{ij}+B_{ij}\sin\theta_{ij}\bigr),\label{eq:pf-th-fund-polar-p}\\
Q_i&=\sum_{j\in\mathcal N}V_{m,i}V_{m,j}
\bigl(G_{ij}\sin\theta_{ij}-B_{ij}\cos\theta_{ij}\bigr).\label{eq:pf-th-fund-polar-q}
\end{align}
\end{corollary}

\begin{proof}
$\overline{Y_{ij}V_j}=(G_{ij}-\mathrm jB_{ij})V_{m,j}e^{-\mathrm j\theta_j}$，
乘以 $V_i=V_{m,i}e^{\mathrm j\theta_i}$ 得
$V_{m,i}V_{m,j}(G_{ij}-\mathrm jB_{ij})e^{\mathrm j\theta_{ij}}$；
展开 $e^{\mathrm j\theta_{ij}}=\cos\theta_{ij}+\mathrm j\sin\theta_{ij}$，
取实部、虚部分别即式~\eqref{eq:pf-th-fund-polar-p}、\eqref{eq:pf-th-fund-polar-q}。
\end{proof}

式~\eqref{eq:pf-th-fund-si}--\eqref{eq:pf-th-fund-polar-q} 是潮流计算的
标准模型~\cite{pff:stagg-elabiad,pff:arrillaga-watson}。其实现层转写（逐式对照
源码、ZIP 负荷注入、DC 侧 $P=V_{dc}\odot(G_{dc}V_{dc})$ 等）见本手册实现层
"装配与 SolverData"章与"统一 Newton 章"，此处不再重复。

\subsection{潮流问题的数学分类}\label{sec:pf-th-fund-classification}

\subsubsection{节点类型}

\begin{definition}[PQ、PV、V$\theta$ 节点]\label{def:pf-th-fund-bus-types}
按节点的已知量/未知量划分~\cite{pff:stott-review}：
\begin{itemize}
  \item \textbf{PQ 节点}：指定 $P_i^{\mathrm{spec}},Q_i^{\mathrm{spec}}$；
        未知量 $\theta_i,V_{m,i}$；对应负荷母线与无定压控制的电源母线；
  \item \textbf{PV 节点}：指定 $P_i^{\mathrm{spec}},V_{m,i}^{\mathrm{spec}}$；
        未知量 $\theta_i$（$Q_i$ 由解出后回算并受 $[Q_i^{\min},Q_i^{\max}]$ 校核）；
        对应具备自动电压调节的发电机母线；
  \item \textbf{V$\theta$ 节点（slack/参考节点）}：指定 $\theta_i$（通常 $0$）
        与 $V_{m,i}$；$P_i,Q_i$ 完全自由，由功率平衡残差吸收。
\end{itemize}
\end{definition}

\begin{proposition}[方程-未知数计数]\label{prop:pf-th-fund-counting}
$n$ 节点系统含 $n_{pv}$ 个 PV 节点、$1$ 个 V$\theta$ 节点，其余为 PQ 节点，
则潮流问题的未知量个数与独立方程个数均为
\begin{equation}
m=2n-2-n_{pv}.
\end{equation}
\end{proposition}

\begin{proof}
未知量：角度 $\theta_i$ 对全部非参考节点自由，共 $n-1$ 个；电压幅值 $V_{m,i}$
仅对 PQ 节点自由，共 $n-1-n_{pv}$ 个；合计 $2n-2-n_{pv}$。方程：有功平衡
$P_i^{\mathrm{spec}}-P_i^{\mathrm{calc}}=0$ 对全部非参考节点列出，共 $n-1$ 个；
无功平衡仅对 PQ 节点列出（PV 节点的 $Q$ 不平衡由发电机无功吸收、V$\theta$
节点的 $P,Q$ 均自由），共 $n-1-n_{pv}$ 个；合计 $2n-2-n_{pv}$。两者相等，
问题为方阵非线性方程组 $F(x)=0$，$F:\mathbb R^m\to\mathbb R^m$。
\end{proof}

\begin{remark}[参考节点的双重必要性]\label{rem:pf-th-fund-slack}
V$\theta$ 节点的设置有两个独立的数学理由：
(i) \emph{角度规范}：式~\eqref{eq:pf-th-fund-polar-p}--\eqref{eq:pf-th-fund-polar-q}
中 $F$ 仅通过差 $\theta_{ij}$ 依赖于角度，全体角度同时平移不改变任何失配——
角度空间带有一个 $S^1$ 对称性，不固定一个参考角则雅可比恒有一维零空间；
(ii) \emph{损耗悬空}：网络损耗在求解之前未知，$n$ 个节点的有功指定值不能
同时精确满足，必须有一个节点的 $P$（与 $Q$）退出方程组以吸收差额。
多岛系统每个连通 AC 岛各需一个参考节点；把损耗分摊到多台机组的分布式松弛
是上述单点吸收的推广，不改变计数结论。
\end{remark}

\subsubsection{解的存在性与多解性}

潮流方程是关于 $(\theta,V_m)$ 的二次型代数方程组（固定 $V_m$ 时对
$e^{\mathrm j\theta}$ 为双线性）。其解集一般不是单点：

\begin{example}[两节点系统的解析解与鼻点]\label{ex:pf-th-fund-twobus}
节点 1 为参考节点 $E\angle0$，节点 2 经纯电抗 $X$ 线路受电，负荷
$P_L+\mathrm jQ_L$。记 $V=V_{m,2}$、$\delta=-\theta_2$，则
\begin{equation}
P_L=\frac{EV}{X}\sin\delta,\qquad
Q_L=\frac{EV}{X}\cos\delta-\frac{V^2}{X}.
\end{equation}
消去 $\delta$（利用 $\sin^2+\cos^2=1$）得关于 $V^2$ 的二次方程
\begin{equation}
V^4+\bigl(2Q_LX-E^2\bigr)V^2+X^2\bigl(P_L^2+Q_L^2\bigr)=0,
\label{eq:pf-th-fund-v4}
\end{equation}
\begin{equation}
V^2=\frac{E^2-2Q_LX\pm\sqrt{\Delta}}{2},\qquad
\Delta=E^4-4E^2Q_LX-4X^2P_L^2 .
\label{eq:pf-th-fund-v2}
\end{equation}
$\Delta>0$ 且 $E^2>2Q_LX$ 时存在\textbf{两个}正根：高电压根（工程可行解）与
低电压根；$\Delta=0$ 给出鞍结分岔（鼻点）$P_L^{\max}=\sqrt{E^4-4E^2Q_LX}/(2X)$；
$\Delta<0$ 时无实解。数值上取 $E=1$、$X=0.5$、$Q_L=0$：$P_L=0.5$ 时
$V^2=(1\pm\sqrt{0.75})/2$，即 $V\approx0.966$ 或 $V\approx0.259$；
$P_L=1$ 时两根于 $V=1/\sqrt2$ 处汇合（鼻点）；$P_L>1$ 无解。
\end{example}

\begin{remark}[一般网络的多解性与不可解性]\label{rem:pf-th-fund-multiplicity}
例~\ref{ex:pf-th-fund-twobus} 的定性结论对一般网络成立：潮流方程的实解个数
通常随规模增长，且重负荷参数下可\textbf{不存在}实解（鼻点之外）。两点推论：
(i) Newton 法不收敛\textbf{不能}作为问题无解的证明——它只说明初值不在任何
解的局部收敛域内被给定迭代法捕获；(ii) 工程上关心的是高电压物理解，多解的
枚举/全解搜索属于超出常规潮流的范畴。
\end{remark}

\begin{gapnote}
平台提供 HELM（全纯嵌入）求解器
（\srcpath{src/power\_flow/helm\_solver.cpp:HelmSolver::solve}），
其方法论是沿解析延拓收敛到与平启动连通的物理解，而非枚举多解；
全解枚举（如消元法、区间分析）未实现，也无计划对应接口。
\end{gapnote}

\subsubsection{无功限值与活动集表述}

PV 节点的无功限值使潮流问题严格地说不再是纯方程组，而是带互补约束的
混合问题：
\begin{equation}
Q_g^{\min}\le Q_g(x)\le Q_g^{\max},\qquad
\text{限内时 }V_m=V^{\mathrm{set}},\ \text{触限时 }Q_g=Q_g^{\pm},
\label{eq:pf-th-fund-qlimit}
\end{equation}
三态（下限/限内/上限）可用中值互补函数
$\mathrm{mid}(V_m-V^{\mathrm{set}},\,Q_g-Q_g^{\min},\,Q_g-Q_g^{\max})=0$
统一表示~\cite{pff:hintermuller}。工程实现的通行做法是\emph{活动集}策略：
固定当前 PV/PQ 归属 $A$，解纯方程组 $F_A(x)=0$ 至收敛，在收敛点检验隐含
无功，越限则批量切换为 PQ 并钉在限值，重复至无越限；恢复侧设迟滞与单次
恢复审计以防止活动集循环。平台的活动集契约（含固定超集雅可比布局与
CHKS 光滑化的半光滑 Newton 替代）全文见
\srcpath{docs/theory/pv\_pq\_switching\_contract.md}，实现层转写见本手册
"统一 Newton 章"与"鲁棒求解器章"。

\subsection{潮流可解性的充分条件}\label{sec:pf-th-fund-solvability}

第~\ref{sec:pf-th-fund-classification}~节从两母线鼻点（例~\ref{ex:pf-th-fund-twobus}）
定性说明了潮流解可以不存在。本节给出一个\emph{定量}充分条件：当负荷相对短路容量足够
小时，存在唯一的高电压解，且不动点迭代收敛到它。设网络含单一 slack（母线 $0$，
电压 $u_0$）与 PQ 母线集 $L=\{1,\dots,n\}$，$Y_{bus}$ 按 $\{0\}$ 与 $L$ 分块，
子块 $Y_{LL}$ 非奇异（连通接地网络成立）。

\begin{proposition}[PQ 潮流的不动点形式]\label{prop:pf-th-fund-fixedpoint}
记 PQ 母线注入复功率 $s\in\mathbb C^n$（发电为正）。定义\emph{空载电压}
$w:=-Y_{LL}^{-1}Y_{L0}u_0$（$s=0$ 时的解）。则 PQ 潮流方程等价于不动点方程
\begin{equation}
u=T(u):=w+Y_{LL}^{-1}\,\overline{s\oslash u},
\label{eq:pf-th-fund-fixedpoint}
\end{equation}
其中 $\oslash$ 为逐分量除、$\overline{(\cdot)}$ 为复共轭。
\end{proposition}
\begin{proof}
PQ 母线电流注入 $i_L=Y_{LL}u+Y_{L0}u_0$；由 $s_i=u_i\overline{i_i}$ 得
$i_i=\overline{s_i/u_i}$，即 $i_L=\overline{s\oslash u}$。两式相等，左乘 $Y_{LL}^{-1}$、
移项即式~\eqref{eq:pf-th-fund-fixedpoint}。
\end{proof}

\begin{theorem}[压缩可解性条件]\label{thm:pf-th-fund-solvability}
记 $\eta:=\lVert Y_{LL}^{-1}\rVert_\infty$（行和范数）、$\sigma:=\max_i|s_i|$、
$w_{\min}:=\min_i|w_i|>0$。若
\begin{equation}
4\,\eta\,\sigma<w_{\min}^2,
\label{eq:pf-th-fund-solvcond}
\end{equation}
则映射 $T$ 在闭球 $\mathcal B=\{u:\lVert u-w\rVert_\infty\le w_{\min}/2\}$ 上是压缩自映射，
潮流方程~\eqref{eq:pf-th-fund-fixedpoint} 在 $\mathcal B$ 内有\emph{唯一}解，且 Picard
迭代 $u^{(k+1)}=T(u^{(k)})$ 从 $u^{(0)}=w$ 出发几何收敛到它。
\end{theorem}
\begin{proof}
在 $\mathcal B$ 上 $|u_i|\ge|w_i|-\tfrac{w_{\min}}2\ge\tfrac{w_{\min}}2=:\rho$。
\textbf{自映射}：对 $u\in\mathcal B$，
\begin{equation*}
\lVert T(u)-w\rVert_\infty=\bigl\lVert Y_{LL}^{-1}\overline{s\oslash u}\bigr\rVert_\infty
\le\eta\,\max_i\frac{|s_i|}{|u_i|}\le\frac{\eta\sigma}{\rho}
=\frac{2\eta\sigma}{w_{\min}}<\frac{w_{\min}}2,
\end{equation*}
末步用式~\eqref{eq:pf-th-fund-solvcond}，故 $T(u)\in\mathcal B$。
\textbf{压缩}：对 $u,u'\in\mathcal B$，逐分量
$\bigl|\overline{s_i/u_i}-\overline{s_i/u_i'}\bigr|
=|s_i|\,|u_i-u_i'|/(|u_i||u_i'|)\le(\sigma/\rho^2)\lVert u-u'\rVert_\infty$，故
\begin{equation*}
\lVert T(u)-T(u')\rVert_\infty\le\eta\,\frac{\sigma}{\rho^2}\lVert u-u'\rVert_\infty
=\frac{4\eta\sigma}{w_{\min}^2}\lVert u-u'\rVert_\infty,
\end{equation*}
压缩系数 $L=4\eta\sigma/w_{\min}^2<1$。$\mathcal B$ 是完备度量空间中的闭集，Banach
不动点定理给出唯一不动点与几何收敛。
\end{proof}

\begin{remark}[短路容量裕度的读法与两母线校核]\label{rem:pf-th-fund-solvability}
$\eta=\lVert Y_{LL}^{-1}\rVert_\infty$ 是“路径阻抗”量级、$\eta^{-1}$ 度量短路容量，
$w_{\min}$ 为最弱母线的空载电压；条件~\eqref{eq:pf-th-fund-solvcond} 即“负荷 $\times$
路径阻抗 $\ll$ 空载电压$^2$”的短路容量裕度，与前推回代的压缩条件
（式~\eqref{sec:pf-th-cls-eq:bfsc}，第~\ref{sec:pf-th-cls}~章）同源。以两母线系统
（例~\ref{ex:pf-th-fund-twobus}，$Q=0$）校核：$Y_{LL}=1/(\mathrm jX)$、$w=E$、
$\eta=X$、$\sigma=P$，条件给出 $P<E^2/(4X)$，恰为真鼻点 $P_{\max}=E^2/(2X)$ 的一半——
即该充分条件保守约 $2$ 倍，符合“压缩界不触鼻点极限”的一般预期。更锐的（计及负荷
方向与网络结构的）可解性判据见 Bolognani--Zampieri~\cite{pff:bolognani} 与
Simpson-Porco--D\"orfler--Bullo~\cite{pff:simpsonporco}。
\end{remark}

\begin{gapnote}
条件~\eqref{eq:pf-th-fund-solvcond} 作为求解前的可解性预判在平台中\textbf{未实现}：
求解器直接尝试 Newton/BFS，以收敛与否与残差水平事后判定；HELM
（\srcpath{src/power\_flow/helm\_solver.cpp:HelmSolver::solve}）沿解析延拓返回与平启动
连通的物理解，同样不输出式~\eqref{eq:pf-th-fund-solvcond} 型的先验证书。
\end{gapnote}

\subsection{Newton--Raphson 一般理论}\label{sec:pf-th-fund-newton}

\subsubsection{迭代格式与局部二次收敛}

\begin{definition}[Newton 迭代]\label{def:pf-th-fund-newton}
对 $F:\mathbb R^m\to\mathbb R^m$、$F\in C^1$，Newton 迭代为
\begin{equation}
J(x_k)\,\Delta x_k=-F(x_k),\qquad x_{k+1}=x_k+\Delta x_k .
\label{eq:pf-th-fund-newton-step}
\end{equation}
潮流语境下 $F$ 取失配 $F^{\mathrm{spec}}-F^{\mathrm{calc}}$（符号约定不影响
理论），雅可比在极坐标下分块为
$J=\begin{bmatrix}H&N\\ M&L\end{bmatrix}$，
$H=\partial P/\partial\theta$、$N=\partial P/\partial V_m$、
$M=\partial Q/\partial\theta$、$L=\partial Q/\partial V_m$
~\cite{pff:tinney-hart}。
\end{definition}

\begin{lemma}[Banach 扰动引理]\label{lem:pf-th-fund-banach}
设 $A\in\mathbb R^{m\times m}$ 非奇异，$\lVert A^{-1}\rVert\le\beta$，
$E\in\mathbb R^{m\times m}$ 满足 $\beta\lVert E\rVert<1$，则 $A+E$ 非奇异且
\begin{equation}
\bigl\lVert(A+E)^{-1}\bigr\rVert\le\frac{\beta}{1-\beta\lVert E\rVert}.
\end{equation}
\end{lemma}

\begin{proof}
写 $A+E=A(I+A^{-1}E)$。$\lVert A^{-1}E\rVert\le\beta\lVert E\rVert<1$，
Neumann 级数 $\sum_{k\ge0}(-A^{-1}E)^k$ 绝对收敛，其和为 $(I+A^{-1}E)^{-1}$，
且 $\lVert(I+A^{-1}E)^{-1}\rVert\le\sum_k(\beta\lVert E\rVert)^k
=1/(1-\beta\lVert E\rVert)$。于是 $(A+E)^{-1}=(I+A^{-1}E)^{-1}A^{-1}$，
取范数即得结论。
\end{proof}

\begin{theorem}[Newton 法的局部二次收敛]\label{thm:pf-th-fund-quadratic}
设 $F\in C^1$，$J=DF$ 在 $x^*$ 的邻域内 Lipschitz 连续：
$\lVert J(x)-J(y)\rVert\le L\lVert x-y\rVert$；
$F(x^*)=0$ 且 $J(x^*)$ 非奇异，$\lVert J(x^*)^{-1}\rVert=\beta$。
则存在 $\delta>0$，使任意 $x_0\in B(x^*,\delta)$ 出发的 Newton 序列
良定义、收敛到 $x^*$，且
\begin{equation}
\lVert e_{k+1}\rVert\le\beta L\lVert e_k\rVert^2,
\qquad e_k=x_k-x^* .
\label{eq:pf-th-fund-qconv}
\end{equation}
\end{theorem}

\begin{proof}
取 $\delta\le\min\{r,\,1/(2\beta L)\}$，其中 $r$ 为 Lipschitz 邻域半径。
归纳假设 $x_k\in B(x^*,\delta)$：则
$\lVert J(x_k)-J(x^*)\rVert\le L\delta\le1/(2\beta)$，由
引理~\ref{lem:pf-th-fund-banach}（$\varepsilon=1/(2\beta)$，$\beta\varepsilon=1/2<1$），
$J(x_k)$ 非奇异且 $\lVert J(x_k)^{-1}\rVert\le2\beta$，迭代良定义。
由 $F(x_k)-F(x^*)=\int_0^1J(x^*+te_k)e_k\,\mathrm dt$ 与 Newton 方程
$J(x_k)(x_k+\Delta x_k-x^*)=J(x_k)e_k-F(x_k)$，得误差恒等式
\begin{equation}
J(x_k)\,e_{k+1}
=\int_0^1\bigl[J(x_k)-J(x^*+te_k)\bigr]e_k\,\mathrm dt .
\end{equation}
被积函数范数 $\le L(1-t)\lVert e_k\rVert^2$，积分得
$\lVert J(x_k)e_{k+1}\rVert\le\frac{L}{2}\lVert e_k\rVert^2$，
左乘 $J(x_k)^{-1}$ 并用 $2\beta$ 界即得式~\eqref{eq:pf-th-fund-qconv}。
又因 $\beta L\lVert e_k\rVert\le\beta L\delta\le1/2$，
$\lVert e_{k+1}\rVert\le\frac12\lVert e_k\rVert$，序列保持于球内并收缩，
故收敛且为 Q-二次。
\end{proof}

\begin{remark}[潮流场景的条件核对]\label{rem:pf-th-fund-hypotheses}
定理~\ref{thm:pf-th-fund-quadratic} 的三项假设在潮流中分别对应：
(i) $C^1$+Lipschitz：式~\eqref{eq:pf-th-fund-polar-p}--\eqref{eq:pf-th-fund-polar-q}
为三角多项式，任意紧集上自动满足；(ii) $F(x^*)=0$：物理解存在性，见
第~\ref{sec:pf-th-fund-classification}~节，非平凡；(iii) $J(x^*)$ 非奇异：
在鼻点处\textbf{系统性失效}（命题~\ref{prop:pf-th-fund-fold}）。"常规工况、
平启动、$3$--$7$ 次迭代收敛"的工程经验正是二次收敛阶段
$\lVert e_{k+1}\rVert\approx c\lVert e_k\rVert^2$ 的表现：
残差小数位数近似逐次翻倍。
\end{remark}

\begin{remark}[仿射协变性与实现中的缩放]\label{rem:pf-th-fund-affine}
Newton 序列在两组变换下不变：未知量的仿射重参数化 $\tilde x=Ax+b$ 与
方程的左乘非奇异阵 $\tilde F=BF$（精确算术下逐点一致）。这一\emph{仿射协变性}
意味着：对角缩放 $D_F=\mathrm{diag}(1/\max(|F_i^{\mathrm{spec}}|,1))$、
$D_x=\mathrm{diag}(\pi,V_{m,i},\dots)$ 不改变理论迭代轨线，只改变浮点误差
的分布——其实现动机（量纲均衡、改善 $J$ 的行/列范数比）见
第~\ref{sec:pf-th-fund-conditioning}~节。反之，步长截断
（$|\Delta\theta|\le1.5$~rad 等）与线搜索回退\textbf{破坏}纯 Newton 序列，
定理~\ref{thm:pf-th-fund-quadratic} 只在进入全步接受阶段后适用。
\end{remark}

\subsubsection{收敛域与全局化}

定理~\ref{thm:pf-th-fund-quadratic} 只保证\emph{存在}收敛球，未给出其大小。
Newton 迭代的全局吸引域结构复杂（多解系统的吸引域边界可为分形），工程上
以全局化策略替代收敛域刻画~\cite{pff:nocedal-wright}：在 merit 函数
$\phi(x)=\tfrac12\lVert F(x)\rVert^2$ 上做（非单调）Armijo 线搜索、
信赖域（dogleg）、Levenberg--Marquardt 正则化
$(J^{\mathsf T}J+\lambda I)\Delta x=-J^{\mathsf T}F$，或伪瞬态延拓
$(J+I/\delta t)\Delta x=-F$。这些策略把局部方法改造为对更大初值集
"收敛或明确失败"的方法，但都不提供"解存在"的证书。

\begin{gapnote}
定量收敛判据——Kantorovich 定理（以
$h=\beta L\lVert\Delta x_0\rVert\le1/2$ 在首步数据上判定收敛并给出误差界）
与 Smale 的 $\alpha$-理论——在当前实现中\textbf{未}接线：代码不计算
$h$ 型证书，失败回退链（线搜索 $\to$ 对角正则化 $\to$ LM $\to$ PTC，
实现层见"统一 Newton 章"）是纯启发式全局化。
\end{gapnote}

\subsubsection{雅可比奇异的物理含义}

\begin{proposition}[鞍结折点处雅可比必奇异]\label{prop:pf-th-fund-fold}
设单参数潮流 $F(x,\lambda)=0$（$\lambda$ 为负荷/发电水平参数）存在光滑解弧
$(x(s),\lambda(s))$，$s$ 为弧长参数。若在 $s^*$ 处 $\lambda'(s^*)=0$ 且
$x'(s^*)\neq0$（解弧在 $\lambda$ 方向回折），则 $J^*=D_xF(x(s^*),\lambda(s^*))$
奇异，且 $x'(s^*)\in\ker J^*$。
\end{proposition}

\begin{proof}
对恒等式 $F(x(s),\lambda(s))\equiv0$ 关于 $s$ 微分：
\begin{equation}
D_xF\,x'(s)+D_\lambda F\,\lambda'(s)=0 .
\end{equation}
代入 $s=s^*$：$\lambda'(s^*)=0$ 给出 $J^*x'(s^*)=0$，而 $x'(s^*)\neq0$，
故 $J^*$ 有非平凡零空间。
\end{proof}

\begin{remark}[奇异的两类成因与工程区分]\label{rem:pf-th-fund-singularity}
(i) \emph{参数性奇异}：命题~\ref{prop:pf-th-fund-fold} 的折点，物理含义是
鼻点/静态电压稳定极限——随 $\lambda$ 逼近 $\lambda^*$，
$\sigma_{\min}(J)\to0$，$\kappa(J)\to\infty$（见
第~\ref{sec:pf-th-fund-conditioning}~节）；
(ii) \emph{结构性奇异}：无参考节点的孤岛、悬空节点形成的空行/空列、
理想回路导致的秩亏——与负荷水平无关，是建模/拓扑缺陷。两者在线性求解阶段
都表现为分解失败或巨大修正量，但处置完全不同：前者应转入连续潮流定位
极限点~\cite{pff:ajjarapu-christy}，后者应在装配期诊断剔除。实现层以
空行/空列扫描与方程闭合校核处置后者（见"与实现的对应"表）。
\end{remark}
\subsection{直角坐标表述、二次结构与最优乘子法}\label{sec:pf-th-fund-rect}

极坐标式~\eqref{eq:pf-th-fund-polar-p}--\eqref{eq:pf-th-fund-polar-q} 是平台采用的
标准形式，但改用\emph{直角坐标}可揭示潮流失配的一个精确二次代数结构，并由此导出两个
经典结论：电流注入雅可比的常量非对角块，以及沿 Newton 方向的\emph{闭式最优步长}。

\subsubsection{直角坐标功率方程}
\begin{definition}[直角坐标电压与注入电流]\label{def:pf-th-fund-rect}
记 $V_i=e_i+\mathrm jf_i$。注入电流 $I_i=\sum_jY_{ij}V_j$ 的实/虚部
\begin{equation}
a_i:=\operatorname{Re}I_i=\sum_j\bigl(G_{ij}e_j-B_{ij}f_j\bigr),\qquad
b_i:=\operatorname{Im}I_i=\sum_j\bigl(G_{ij}f_j+B_{ij}e_j\bigr)
\end{equation}
均为状态 $x=(e_1,\dots,e_n,f_1,\dots,f_n)$ 的\emph{线性}函数。
\end{definition}

\begin{proposition}[直角坐标注入功率]\label{prop:pf-th-fund-rectpq}
\begin{equation}
P_i=e_ia_i+f_ib_i,\qquad Q_i=f_ia_i-e_ib_i,
\label{eq:pf-th-fund-rectpq}
\end{equation}
均为 $x$ 的齐二次型；母线电压幅值平方 $V_i^2=e_i^2+f_i^2$ 亦为齐二次型。
\end{proposition}
\begin{proof}
$S_i=V_i\overline{I_i}=(e_i+\mathrm jf_i)(a_i-\mathrm jb_i)
=(e_ia_i+f_ib_i)+\mathrm j(f_ia_i-e_ib_i)$，取实、虚部即得。$a_i,b_i$ 线性于 $x$，
故 $P_i,Q_i$ 为 $x$ 的二次齐次多项式。
\end{proof}

\subsubsection{失配映射的精确二次结构}
\begin{theorem}[潮流失配是二次多项式映射]\label{thm:pf-th-fund-quadmap}
以直角坐标未知量 $x$（slack 的 $e,f$ 固定）装配失配 $F(x)$——PQ 母线取
$(\Delta P_i,\Delta Q_i)$、PV 母线取 $(\Delta P_i,\ (V_i^{\mathrm{set}})^2-e_i^2-f_i^2)$——
则 $F$ 是 $x$ 的二次多项式映射：存在常向量 $c_0$、常矩阵 $J_0$ 与\emph{常}对称双线性
形 $\mathcal B[\cdot,\cdot]$ 使
\begin{equation}
F(x)=c_0+J_0x+\tfrac12\mathcal B[x,x],\qquad D^3F\equiv0 .
\label{eq:pf-th-fund-quadmap}
\end{equation}
从而对任意基点 $x_k$ 与方向 $\Delta$，
\begin{equation}
F(x_k+\alpha\Delta)=F(x_k)+\alpha\,J(x_k)\Delta+\tfrac12\alpha^2\,\mathcal B[\Delta,\Delta]
\label{eq:pf-th-fund-exactstep}
\end{equation}
是\emph{精确}恒等式（无截断），其中 $J(x_k)=J_0+\mathcal B[x_k,\cdot]$。
\end{theorem}
\begin{proof}
由命题~\ref{prop:pf-th-fund-rectpq}，$F$ 的每个分量是 $x$ 的二次多项式（常数项、
线性项与二次项之和），故其三阶及以上导数恒为零，Taylor 展开在二阶精确终止，
即式~\eqref{eq:pf-th-fund-quadmap}。把 $x=x_k+\alpha\Delta$ 代入，二次映射只产生
$\alpha^0,\alpha^1,\alpha^2$ 三项，系数分别为 $F(x_k)$、$J(x_k)\Delta$、
$\tfrac12\mathcal B[\Delta,\Delta]$，即式~\eqref{eq:pf-th-fund-exactstep}。
\end{proof}

\begin{remark}[电流注入雅可比的常量非对角块]\label{rem:pf-th-fund-ci}
以\emph{电流失配} $\Delta I_i=\overline{S_i^{\mathrm{spec}}/V_i}-\sum_jY_{ij}V_j$
（实/虚部堆叠）为方程的电流注入 Newton 法~\cite{pff:dacosta} 中，网络项
$\sum_jY_{ij}V_j$ 对 $V_j$ 线性，故非对角块
$\partial(\Delta I_i^{\mathrm{re}},\Delta I_i^{\mathrm{im}})/\partial(e_j,f_j)$
（$i\ne j$）恒等于\emph{常量} $2\times2$ 实块
$\left[\begin{smallmatrix}-G_{ij}&\phantom{-}B_{ij}\\-B_{ij}&-G_{ij}\end{smallmatrix}\right]$，
与运行点无关，只有 $2\times2$ 对角块随本地负荷电流的状态依赖而变。这一“非对角块免
重算”性质是电流注入法单步更省的结构来源。
\end{remark}

\subsubsection{Iwamoto 最优乘子法}
式~\eqref{eq:pf-th-fund-exactstep} 的精确性使沿 Newton 方向的 merit 有\emph{闭式}
最小点。设 $g:=F(x_k)$，Newton 方向 $\Delta_N$ 满足 $J(x_k)\Delta_N=-g$，
记 $c:=\mathcal B[\Delta_N,\Delta_N]$，则由式~\eqref{eq:pf-th-fund-exactstep}，
$F(x_k+\alpha\Delta_N)=(1-\alpha)g+\tfrac12\alpha^2c$。

\begin{proposition}[最优乘子的三次方程]\label{prop:pf-th-fund-iwamoto}
记 $a_0=g^{\mathsf T}g$、$a_1=g^{\mathsf T}c$、$a_2=c^{\mathsf T}c$，则 merit
$\phi(\alpha)=\tfrac12\lVert F(x_k+\alpha\Delta_N)\rVert^2$ 是 $\alpha$ 的\emph{四次}
多项式，其驻点由\emph{三次}方程
\begin{equation}
a_2\,\alpha^3-3a_1\,\alpha^2+2(a_0+a_1)\,\alpha-2a_0=0
\label{eq:pf-th-fund-iwamoto}
\end{equation}
给出。落在 $(0,1]$ 内使 $\phi$ 最小的实根即\emph{最优乘子} $\mu_k$，可由 Cardano
公式闭式求得，代价可忽略~\cite{pff:iwamoto}。
\end{proposition}
\begin{proof}
$\phi(\alpha)=\tfrac12\bigl[(1-\alpha)^2a_0+(1-\alpha)\alpha^2a_1+\tfrac14\alpha^4a_2\bigr]$。
求导 $\phi'(\alpha)=\tfrac12\bigl[-2a_0(1-\alpha)+a_1(2\alpha-3\alpha^2)+a_2\alpha^3\bigr]$，
令其为零并整理即式~\eqref{eq:pf-th-fund-iwamoto}。又 $\phi'(0)=-a_0<0$
（$a_0=\lVert g\rVert^2>0$），保证在 $\alpha>0$ 侧存在使 merit 下降的根。
\end{proof}

\begin{remark}[最优乘子的全局化意义与渐近还原]\label{rem:pf-th-fund-iwamoto}
$\mu_k$ 给出沿 Newton 方向的\emph{精确}最优 merit 下降，无需回溯试探，历史上专用于
病态/发散边缘系统的稳健求解~\cite{pff:iwamoto}。近解处 $\Delta_N\to0$，
$c=\mathcal B[\Delta_N,\Delta_N]$ 为二阶小量，$a_2=\lVert c\rVert^2\to0$，
式~\eqref{eq:pf-th-fund-iwamoto} 退化为 $2(a_0+a_1)\alpha-2a_0\approx0$，
$\mu_k\to1$：方法自动还原为全步 Newton，定理~\ref{thm:pf-th-fund-quadratic} 的
二次收敛得以保留。这与极坐标下 Armijo 线搜索近解 $\alpha\to1$ 殊途同归，差别在于
二次结构给出了\emph{无试探}的闭式步长。
\end{remark}

\begin{gapnote}
平台统一 Newton 求解器采用\textbf{极坐标}功率失配
（\srcpath{src/power\_flow/ac\_kernel.cpp:evaluate\_ac\_kernel\_serial}），全局化用
Armijo/非单调线搜索、LM 与 PTC（见第~\ref{sec:pf-th-glob}~章），
\textbf{未实现}直角坐标 Newton、电流注入 Newton 与 Iwamoto 最优乘子；
式~\eqref{eq:pf-th-fund-iwamoto} 依赖的精确二次结构在极坐标下不成立（三角项使
$D^3F\ne0$）。平台内唯一的电流注入构造是三相 Newton 的定点辅助
（\srcpath{src/power\_flow/three\_phase\_nr.cpp:evaluate\_fixed\_point\_current\_injections}），
非本节的电流失配 Newton 主迭代。
\end{gapnote}
\subsection{病态与数值性态}\label{sec:pf-th-fund-conditioning}

\subsubsection{条件数与扰动界}

\begin{definition}[条件数与向后误差]\label{def:pf-th-fund-cond}
线性系统 $Jx=b$ 的条件数 $\kappa(J)=\lVert J\rVert\lVert J^{-1}\rVert$。
称 $\hat x$ 具有向后误差 $\eta$，若 $\eta$ 是使
$(J+\Delta J)\hat x=b+\Delta b$、$\lVert\Delta J\rVert\le\eta\lVert J\rVert$、
$\lVert\Delta b\rVert\le\eta\lVert b\rVert$ 成立的最小相对扰动。
\end{definition}

\begin{theorem}[线性系统扰动界]\label{thm:pf-th-fund-perturb}
设 $Jx=b$、$(J+\Delta J)(x+\Delta x)=b+\Delta b$，记
$\kappa=\kappa(J)$、$\rho=\lVert\Delta J\rVert/\lVert J\rVert$，
且 $\kappa\rho<1$，则
\begin{equation}
\frac{\lVert\Delta x\rVert}{\lVert x\rVert}
\le\frac{\kappa}{1-\kappa\rho}
\Bigl(\frac{\lVert\Delta b\rVert}{\lVert b\rVert}+\rho\Bigr).
\label{eq:pf-th-fund-perturb}
\end{equation}
\end{theorem}

\begin{proof}
两式相减：$J\Delta x=\Delta b-\Delta J\,x-\Delta J\,\Delta x$。
$J+\Delta J$ 非奇异（引理~\ref{lem:pf-th-fund-banach}，
$\lVert J^{-1}\rVert\lVert\Delta J\rVert\le\kappa\rho<1$）。左乘 $J^{-1}$
并取范数：
\begin{equation}
\lVert\Delta x\rVert
\le\lVert J^{-1}\rVert\bigl(\lVert\Delta b\rVert+\lVert\Delta J\rVert\,
\lVert x\rVert\bigr)+\lVert J^{-1}\rVert\lVert\Delta J\rVert\lVert\Delta x\rVert .
\end{equation}
移项除以 $1-\kappa\rho$，再用
$\lVert\Delta b\rVert/\lVert x\rVert
=\lVert J\rVert\,\lVert\Delta b\rVert/(\lVert J\rVert\lVert x\rVert)
\le\lVert J\rVert\,\lVert\Delta b\rVert/\lVert b\rVert$
即得式~\eqref{eq:pf-th-fund-perturb}。
\end{proof}

\begin{remark}[浮点含义与潮流量级]\label{rem:pf-th-fund-float}
部分主元 LU 的向后误差 $\eta\approx u\cdot\rho_{\mathrm{growth}}$
（增长因子，电力雅可比上实践接近 $1$）~\cite{pff:golub-vanloan}。
由式~\eqref{eq:pf-th-fund-perturb}，每步 Newton 修正的相对前向误差量级为
$u\,\kappa(J)$：$\kappa(J)\sim10^{8}$ 时双精度下修正量只剩约 8 位有效数字，
残差无法压到 $10^{-8}$ 以下；$\kappa(J)\sim10^{16}$ 时修正量完全失真。
这给出收敛容差 \fld{tol} 的一个理论下界 $\approx u\,\kappa(J)$：
把容差设到病态系统分辨率之下，表现为"残差平台期"而非继续二次收敛。
\end{remark}

\subsubsection{高 \texorpdfstring{$R/X$}{R/X} 比与解耦失效}

快速解耦潮流（FDPF）的推导假设~\cite{pff:stott-alsac}：
(i) $|G_{ij}|\ll|B_{ij}|$（即 $R/X\ll1$）；(ii) $|\theta_{ij}|$ 小，
$\cos\theta_{ij}\approx1$；(iii) $Q_i\ll B_{ii}V_i^2$。三者使
$N=\partial P/\partial V_m$ 与 $M=\partial Q/\partial\theta$ 近似可忽略，
$\theta$ 与 $V_m$ 解耦。配电网络 $R/X$ 常在 $1$ 附近甚至更大，假设 (i) 失效，
交叉块量级与对角块可比，FDPF 迭代矩阵谱半径逼近或超过 $1$，表现为线性收敛
变慢直至发散。Newton 法的渐近二次收敛\textbf{不}依赖解耦假设，但高 $R/X$
使 $P$--$V$ 强耦合、雅可比失去块对角主导，配合辐射状拓扑（无环对角的
"链式"结构）整体条件数上升；此时专用替代算法（前推回代类，利用辐射树结构
逐支路递推）在鲁棒性上优于直接 Newton~\cite{pff:stott-review}。

\subsubsection{重负荷、弱网与电压崩溃}

随负荷参数 $\lambda\to\lambda^*$（鼻点），命题~\ref{prop:pf-th-fund-fold}
给出 $\sigma_{\min}(J)\to0$；由于 $\lVert J\rVert$ 温和变化，
$\kappa(J)\to\infty$。与之平行的物理图景：短路容量低（弱网）的节点
$\partial V/\partial Q$ 灵敏度大，同样放大 $J^{-1}$ 的范数。重负荷/弱网病态
因此\emph{不是}浮点异常，而是临近静态稳定极限的准确数值反映：
雅可比的左零向量（奇异向量）同时是负荷参数变化对失配影响的方向，
这一对应是连续潮流以弧长参数化翻越鼻点的理论基础
~\cite{pff:ajjarapu-christy,pff:vancutsem-vournas,pff:kundur}。

\begin{gapnote}
当前实现未计算 $\kappa(J)$ 或其标准估计（如 Hager--Higham $1$-范数估计器、
$\sigma_{\min}$ 监控）：代码仅计算行列范数代理
\srcpath{src/power\_flow/newton\_solver.cpp:estimate\_condition\_proxy}，
并以代理越阈触发 NK-GMRES 回退（默认关闭）。鼻点附近的奇异向量分析、
按 $\sigma_{\min}$ 的裕度指标均未实现；电压稳定极限定位由连续潮流求解器
（弧长延拓，实现层见"CPF 章"）承担，而非由病态诊断推导。
\end{gapnote}

\subsection{稀疏技术理论}\label{sec:pf-th-fund-sparsity}

\subsubsection{消元图与填充}

潮流雅可比（固定 PV/PQ 活动集下）具有与 $Y_{bus}$ 图同构的对称稀疏模式：
每个 AC 节点至多四个块。对这类矩阵，稀疏 LU 的成本完全由\emph{消元顺序}
决定，其理论语言是消元图~\cite{pff:george-liu,pff:davis-suitesparse}。

\begin{definition}[矩阵图与消元图]\label{def:pf-th-fund-elimgraph}
对称模式矩阵 $A\in\mathbb R^{n\times n}$ 的图 $G(A)=(V,E)$：
$V=\{1,\dots,n\}$，$(i,j)\in E\iff a_{ij}\neq0$（$i\neq j$）。
依次消去顶点 $v_1,v_2,\dots$：消去 $v$ 时把 $v$ 的当前邻居两两连边
（已有边不重复），再删除 $v$；新产生的边称为\textbf{填充}（fill-in），
对应 LU 因子中 $A$ 原本为零位置上的非零元。
\end{definition}

\begin{proposition}[Parter 规则]\label{prop:pf-th-fund-parter}
以顶点 $v$ 为主元的消元步，在消元图中把 $v$ 的邻居集变成团（clique）：
对任意两个邻居 $i,k$，消元后 $(i,k)$ 为边。
\end{proposition}

\begin{proof}
对称高斯消元一步的 Schur 补更新为
$\tilde a_{ik}=a_{ik}-a_{iv}a_{vv}^{-1}a_{vk}$。
$i,k$ 同为 $v$ 的邻居时 $a_{iv}a_{vk}\neq0$，修正项一般非零，故
$\tilde a_{ik}\neq0$ 无论 $a_{ik}$ 是否为零——已有边保留，缺失处产生填充边。
精确相消（数值对消）是零测集事件，符号分析按一般位置处理。
\end{proof}

\begin{example}[星形网络的排序敏感性]\label{ex:pf-th-fund-star}
中心节点 $1$ 连接叶节点 $2,\dots,5$（如主 slack 汇接多条馈线）。
$A$ 的非零元 $13$ 个（$5$ 对角 $+$ $2\times4$ 非对角）。
\begin{itemize}
  \item 顺序 $(1,2,3,4,5)$：消去中心时其邻居为 $\{2,3,4,5\}$，按
        命题~\ref{prop:pf-th-fund-parter} 形成 $4$ 团，产生
        $\binom{4}{2}=6$ 条填充边；$L+U$ 非零元 $=13+2\times6=25$，
        且后续每步面对稠密团，消元工作量 $\Theta(4^2)$；
  \item 顺序 $(2,3,4,5,1)$：消去任一叶片时其唯一邻居是中心，不产生填充；
        最后消去中心。填充为零，$L+U$ 非零元保持 $13$，工作量 $\Theta(n)$。
\end{itemize}
同一矩阵，仅排序不同，填充与工作量相差一个量级——这就是排序理论要
形式化的现象。
\end{example}

\begin{proposition}[填充路径刻画]\label{prop:pf-th-fund-fillpath}
设 $F=L+U$ 为某消元顺序下的填充矩阵。对 $i>j$，$(i,j)\in G(F)$ 当且仅当
在 $G(A)$ 中存在一条从 $i$ 到 $j$ 的路径，其所有中间顶点的消元序号
小于 $j$（即小于两端序号）。
\end{proposition}

\begin{proof}[证明要点]
对消元步数归纳。基础步：原边显然满足（无中间顶点）。归纳步：消去 $v$
时由命题~\ref{prop:pf-th-fund-parter} 新产生的边 $(i,k)$ 有路径
$i$--$v$--$k$，其中 $v$ 先于 $i,k$ 消元；而既有边 $(i,j)\in G(F)$ 若在
此前各步产生，其见证路径的中间顶点均已被消去、序号小于 $j$。反向：
任给满足条件的路径，沿路径上序号最大的中间顶点归纳分裂为两条更短的
合法路径，最终归结为原边或填充边。严格归纳见 George--Liu
~\cite{pff:george-liu}。
\end{proof}

命题~\ref{prop:pf-th-fund-fillpath} 把"哪些位置会填充"完全归结为图论问题：
填充 $=$ 低序号顶点介导的可达性。排序算法的目标即选消元顺序使这种
可达性最小化。

\subsubsection{最小度与近似最小度排序}

\begin{definition}[最小度（MD）排序]\label{def:pf-th-fund-md}
贪心策略：每步在当前消元图中选度（外部度）最小的顶点消元，平局任意打破；
即 Tinney 方案 2~\cite{pff:tinney-walker}。直观依据：消去度 $d$ 的顶点至多
产生 $\binom{d}{2}$ 条填充、$d^2$ 量级算术运算，局部最小化度近似全局
最小化填充。
\end{definition}

\begin{remark}[AMD 的原理]\label{rem:pf-th-fund-amd}
精确 MD 的瓶颈在于消去顶点后邻居集的团更新与外部度重算
（最坏 $O(n\cdot nnz)$）。近似最小度（AMD）~\cite{pff:amestoy-amd} 以
\emph{商图}表示消元过程：已消元顶点被吸收为"元素"，未消元顶点的邻接
以元素集合表示，外部度用上界近似代替精确计算；配合超变量（结构相同的行
合并）与质量吸收，AMD 在实践中达到接近 MD 的填充质量而开销低一个量级以上，
是 SuiteSparse（KLU/UMFPACK 内部）的默认排序。对规则网孔类网络，嵌套分割
（递归几何/图二分）另有严格界：二维 $n\times n$ 网格上填充
$O(n^2\log n)$、分解工作量 $O(n^3)$~\cite{pff:george-nested}；电力网近似
平面图，AMD 的实际表现通常已接近嵌套分割。
\end{remark}

\subsubsection{填充分析与符号/数值分解分离}

由命题~\ref{prop:pf-th-fund-fillpath}，$L+U$ 的非零模式是 $A$ 的模式与
消元顺序的纯函数，与数值无关。因此稀疏直接法天然分两层：
\textbf{符号分解}（排序 + 模式确定，一次）与\textbf{数值分解}
（填充数值，每次 Newton 迭代）。固定 PV/PQ 活动集与网络拓扑时，雅可比
模式逐迭代不变，符号分解可跨迭代、跨求解复用——这是实现层模式缓存
（JacobianPattern/analyzePattern 一次性构建，迭代内仅数值重分解）的
理论依据。PV$\to$PQ 切换改变行/列结构时传统布局须重建模式，
固定超集布局（所有非参考母线预留 $V_m$ 列与 $Q$ 行）以常量填充冗余换取
模式不变性，其契约见 \srcpath{docs/theory/pv\_pq\_switching\_contract.md}。

\begin{gapnote}
平台\textbf{未自研}排序与块三角分解（BTF）：平衡潮流的排序委托给
KLU 的符号分析（其内部默认为 AMD，并可经环境变量切换 COLAMD/嵌套分割，
见 \srcpath{MIPSolvers/src/engine/kernel/linear\_algebra/linear\_solver.cpp:make\_default\_sparse\_solver}）；
三相混合求解器直接使用 Eigen 的 COLAMDOrdering
（\srcpath{src/power\_flow/three\_phase\_hybrid.cpp} 内 SparseLU 配置）。
本节的 MD/AMD/嵌套分割理论用于理解因子填充来源，不对应平台自有代码；
自定义排序接口亦未暴露。
\end{gapnote}

\subsection{数值算例与基准}\label{sec:pf-th-fund-numexp}

本节全部数字由 \srcpath{tools/pf\_doc\_benchmark.cpp}（模式 \texttt{nr}）在 MATPOWER
标准算例（\srcpath{data/case*.m}）上、经平台默认统一 Newton 求解器（极坐标、KLU 稀疏
分解、平启动、$\mathrm{tol}=10^{-10}$）实测生成，可随代码演进重跑复现。迭代次数与残差
轨迹为机器无关量；求解墙钟时间仅供量级参考（Apple~M4，单线程）。

\subsubsection{收敛的规模弱依赖性}
表~\ref{tab:pf-th-fund-nrconv} 给出 $9\sim300$ 母线算例的收敛数据。Newton 维数
$m=2(n-1)$（固定超集布局：非参考母线一律保留 $\theta$ 与 $V_m$ 两列，见
第~\ref{sec:pf-th-fund-sparsity}~节末）。迭代次数稳定在 $4\sim8$ 次、几乎不随规模增长，
最终残差压到 $10^{-13}$ 量级——正是注记~\ref{rem:pf-th-fund-hypotheses} 所述“平启动、
个位数次迭代收敛”的定量印证。

\begin{table}[H]\centering\footnotesize
\caption{统一 Newton 潮流在 MATPOWER 算例上的收敛数据（平启动，$\mathrm{tol}=10^{-10}$）}
\label{tab:pf-th-fund-nrconv}
\begin{tabular}{@{}lrrrrrr@{}}
\toprule
算例 & 母线数 $n$ & Newton 维数 $m$ & 迭代次数 & 最终残差 & 条件数代理 & PV$\to$PQ 切换\\
\midrule
case9   & 9   & 16  & 5 & $1.7\times10^{-14}$ & $5.9\times10^{2}$ & 0\\
case14  & 14  & 26  & 4 & $9.6\times10^{-15}$ & $1.9\times10^{3}$ & 0\\
case30  & 30  & 58  & 5 & $8.9\times10^{-15}$ & $2.0\times10^{4}$ & 0\\
case57  & 57  & 112 & 4 & $3.4\times10^{-12}$ & $2.2\times10^{4}$ & 0\\
case118 & 118 & 234 & 7 & $8.5\times10^{-14}$ & $1.3\times10^{5}$ & 6\\
case300 & 300 & 598 & 8 & $4.3\times10^{-13}$ & $1.3\times10^{9}$ & 10\\
\bottomrule
\end{tabular}
\end{table}

条件数代理列由 \srcpath{estimate\_condition\_proxy}（行/列 $\infty$-范数之积，\emph{非}真
$\kappa_2$）给出，随规模与负荷压力单调上升。case300 的代理达 $10^{9}$ 量级，但其真解
仍收敛到 $4\times10^{-13}$：若真 $\kappa_2$ 果为 $10^{9}$，由注记~\ref{rem:pf-th-fund-float}
的 $u\,\kappa$ 界残差应止步于 $\sim10^{-7}$，故该行/列范数代理\emph{显著高估}真
$\kappa_2$，只宜作趋势指标而非绝对条件数。

\subsubsection{二次收敛的逐迭代验证}
\begin{example}[case30 的二次收敛率]\label{ex:pf-th-fund-quadconv}
下表列出 case30 每次 Newton 迭代记录的最大失配范数 $r_k$（$k=0$ 为平启动初始失配）
与二次收敛率 $r_k/r_{k-1}^2$。比率在 $r_2,r_3$ 步稳定于 $\approx0.22$（有界），
即 $r_{k+1}\approx0.22\,r_k^2$——定理~\ref{thm:pf-th-fund-quadratic} 的
$\lVert e_{k+1}\rVert\le\beta L\lVert e_k\rVert^2$ 在此对应 $\beta L\approx0.2$。
末步 $r_4$ 触及 $\sim10^{-14}$ 的浮点残差地板，比率失真（$r_3^2\approx9\times10^{-19}$
已低于双精度可分辨范围），与注记~\ref{rem:pf-th-fund-float} 的残差平台期吻合。
\begin{center}\footnotesize
\begin{tabular}{@{}crl@{}}
\toprule
迭代 $k$ & 最大失配 $r_k$ & 二次率 $r_k/r_{k-1}^2$\\
\midrule
0 & $3.93\times10^{-1}$ & —\\
1 & $1.63\times10^{-2}$ & $1.06\times10^{-1}$\\
2 & $6.59\times10^{-5}$ & $2.47\times10^{-1}$\\
3 & $9.57\times10^{-10}$ & $2.20\times10^{-1}$\\
4 & $8.88\times10^{-15}$ & （残差地板）\\
\bottomrule
\end{tabular}
\end{center}
\end{example}

\subsubsection{无功限值切换对 Newton 轨迹的扰动}
\begin{example}[case300 的两段式收敛]\label{ex:pf-th-fund-qswitch}
case300 逐迭代记录的最大失配为（竖线标注一次 PV$\to$PQ 切换处）
\begin{center}\footnotesize
$9.27,\ 2.59,\ 3.54{\times}10^{-1},\ 8.45{\times}10^{-3},\ 4.65{\times}10^{-6},\
1.38{\times}10^{-12}\ \big\|\ 3.00{\times}10^{-3},\ 2.04{\times}10^{-7},\
4.26{\times}10^{-13}.$
\end{center}
轨迹呈三个阶段：(i) \emph{大修正区}（$9.27\to2.59\to0.35$，比率非二次），
对应初值远离解、超出定理~\ref{thm:pf-th-fund-quadratic} 的二次收敛邻域，
线搜索与步长限幅等全局化机制主导（第~\ref{sec:pf-th-fund-newton}~节收敛域讨论）；
(ii) \emph{二次收敛尾}（$8.45{\times}10^{-3}\to4.65{\times}10^{-6}\to1.38{\times}10^{-12}$，
残差近似逐步平方，比率 $\approx0.065$）；(iii) 竖线处一次 \textbf{无功限值 PV$\to$PQ 切换}
改变活动集、重新注入 $\sim3\times10^{-3}$ 的失配，之后再次二次收敛到 $4.3\times10^{-13}$。
这段真实轨迹同时印证三件事：远离解的全局化行为、近解的二次收敛，以及
式~\eqref{eq:pf-th-fund-qlimit} 的无功限值活动集切换如何以“重启 Newton”的形式表现——
纯 Newton 序列被活动集变更打断，正是注记~\ref{rem:pf-th-fund-affine} 所述“限值切换
破坏纯 Newton 序列”的实测样貌。
\end{example}

\subsection{与实现的对应}\label{sec:pf-th-fund-implmap}
\begin{center}\footnotesize
\begin{longtable}{@{}L{6.8cm}L{8.6cm}@{}}
\toprule
理论条目 & 实现状态\\
\midrule
\endfirsthead
\toprule
理论条目 & 实现状态\\
\midrule
\endhead
$\pi$ 等值支路组装（串联导纳 + 两端充电 $b_c/2$ + 节点/shunt 并联） &
\implfull{src/power\_flow/admittance\_builder.cpp:build\_admittance\_matrix}\\
复变比 $t=ae^{\mathrm j\varphi}$ 两端口（式~\eqref{eq:pf-th-fund-tport}，
含移相器非对称） &
\implfull{src/power\_flow/admittance\_builder.cpp:build\_admittance\_matrix}（幅值 \fld{tap}、移相角 \fld{shift\_deg} 合成复数 \fld{tap}，$Y_{ft}=-y_s/\overline{t}$、$Y_{tf}=-y_s/t$）\\
节点功率方程极坐标注入（推论~\ref{cor:pf-th-fund-polar}） &
\implfull{src/power\_flow/ac\_kernel.cpp:evaluate\_ac\_kernel\_serial}\\
PQ/PV/V$\theta$ 分类与方程-未知数闭合（命题~\ref{prop:pf-th-fund-counting}） &
\implfull{src/power\_flow/newton\_solver.cpp:build\_jacobian\_context}（闭合计数校核见同文件 scan\_empty\_rows\_cols）\\
多岛各自参考节点（注记~\ref{rem:pf-th-fund-slack}） &
\implfull{src/power\_flow/newton\_solver.cpp:build\_jacobian\_context}（每个孤立 AC 区域各自保留参考；DC 岛参考另有规划函数）\\
无功限值活动集策略（式~\eqref{eq:pf-th-fund-qlimit}） &
\implfull{src/power\_flow/newton\_solver.cpp:solve}（收敛点批量进入 + 迟滞恢复 + 单次恢复审计；契约见 docs/theory/pv\_pq\_switching\_contract.md）\\
中值互补/CHKS 光滑化的半光滑 Newton &
\implfull{src/power\_flow/newton\_solver.cpp:solve}（由 \fld{enable\_semi\_smooth\_newton} 启用，$\mu$ 延拓见 robust\_nonlinear\_options）\\
多解枚举与全解搜索（注记~\ref{rem:pf-th-fund-multiplicity}） &
\implnone（HELM 求解器 helm\_solver.cpp:HelmSolver::solve 仅返回与平启动连通的物理解，非枚举）\\
可解性的压缩充分条件（定理~\ref{thm:pf-th-fund-solvability}，
式~\eqref{eq:pf-th-fund-solvcond}） &
\implnone（无求解前可解性预判；以收敛与否事后判定）\\
Newton 迭代与局部二次收敛（定理~\ref{thm:pf-th-fund-quadratic}） &
\implfull{src/power\_flow/newton\_solver.cpp:solve}（实现层转写见本手册"统一 Newton 章"）\\
Kantorovich/$\alpha$ 定量收敛证书 &
\implnone（以线搜索/信赖域/LM/PTC 启发式全局化回退链替代）\\
折点雅可比奇异与鼻点定位（命题~\ref{prop:pf-th-fund-fold}） &
\implfull{src/power\_flow/voltage\_stability.cpp:solve}（CpfSolver 弧长延拓翻越鼻点）\\
结构性奇异诊断（注记~\ref{rem:pf-th-fund-singularity}） &
\implpart{空行/空列扫描与方程闭合校核已实现（scan\_empty\_rows\_cols）；诊断不区分参数性/结构性成因，也不输出奇异向量}\\
条件数监控与扰动界（定理~\ref{thm:pf-th-fund-perturb}） &
\implpart{仅行列范数代理 estimate\_condition\_proxy 与可选 NK-GMRES 回退；无 Hager--Higham 或 $\sigma_{\min}$ 估计}\\
仿射缩放 $D_F,D_x$（注记~\ref{rem:pf-th-fund-affine}） &
\implfull{src/power\_flow/nonlinear\_scaling.cpp:build\_nonlinear\_scaling}\\
高 $R/X$ 配网的前推回代替代 &
\implfull{src/power\_flow/distribution\_power\_flow.cpp:solve\_distribution\_pf}（辐射状配网 BFS；FDPF 本体见 fdpf\_solver.cpp:FDPFSolver::solve）\\
填充路径/消元图分析与符号-数值分解分离 &
\implfull{src/power\_flow/jacobian\_pattern.cpp:build\_jacobian\_pattern}（模式一次构建、迭代内仅数值累加；符号分析缓存见 newton\_solver.cpp:solve 的 PatternCache）\\
MD/AMD/嵌套分割排序算法 &
\implpart{平台未自研：KLU 默认 AMD（MIPSolvers 端 make\_default\_sparse\_solver），三相混合用 Eigen COLAMDOrdering；无自定义排序接口}\\
解析雅可比的差分校验工具 &
\implfull{src/power\_flow/jacobian\_check.cpp:check\_jacobian}（中心差分 $h=10^{-6}$）\\
直角坐标失配的精确二次结构与步长恒等式（定理~\ref{thm:pf-th-fund-quadmap}，
式~\eqref{eq:pf-th-fund-exactstep}） &
\implnone（主迭代用极坐标功率失配，三角项使 $D^3F\ne0$）\\
电流注入 Newton 的常量非对角块（注记~\ref{rem:pf-th-fund-ci}） &
\implpart{仅三相 Newton 含定点电流注入辅助
three\_phase\_nr.cpp:evaluate\_fixed\_point\_current\_injections；单相/正序主迭代为极坐标功率失配}\\
Iwamoto 最优乘子闭式步长（命题~\ref{prop:pf-th-fund-iwamoto}，
式~\eqref{eq:pf-th-fund-iwamoto}） &
\implnone（全局化改用 Armijo/非单调线搜索、LM、PTC，见第~\ref{sec:pf-th-glob}~章）\\
收敛性数值基准（表~\ref{tab:pf-th-fund-nrconv}、例~\ref{ex:pf-th-fund-quadconv}、
\ref{ex:pf-th-fund-qswitch}） &
\implfull{tools/pf\_doc\_benchmark.cpp}（模式 nr；调用 api/hacdcpf.hpp:solve\_power\_flow）\\
\bottomrule
\end{longtable}
\end{center}

\subsection{参考文献}
{\renewcommand{\section}[2]{}%
\begin{thebibliography}{9}\small
\bibitem{pff:stagg-elabiad} G.~W. Stagg, A.~H. El-Abiad, \emph{Computer Methods
in Power System Analysis}, McGraw-Hill, 1968.
\bibitem{pff:arrillaga-watson} J.~Arrillaga, N.~R. Watson, \emph{Computer
Modelling of Electrical Power Systems}, 2nd ed., Wiley, 2001.
\bibitem{pff:tinney-hart} W.~F. Tinney, C.~E. Hart, Power flow solution by
Newton's method, \emph{IEEE Transactions on Power Apparatus and Systems},
vol. PAS-86, no. 11, pp. 1449--1460, 1967.
\bibitem{pff:stott-alsac} B.~Stott, O.~Alsa\c{c}, Fast decoupled load flow,
\emph{IEEE Transactions on Power Apparatus and Systems}, vol. PAS-93, no. 3,
pp. 859--869, 1974.
\bibitem{pff:stott-review} B.~Stott, Review of load-flow calculation methods,
\emph{Proceedings of the IEEE}, vol. 62, no. 7, pp. 916--929, 1974.
\bibitem{pff:ortega-rheinboldt} J.~M. Ortega, W.~C. Rheinboldt,
\emph{Iterative Solution of Nonlinear Equations in Several Variables},
Academic Press, 1970.
\bibitem{pff:nocedal-wright} J.~Nocedal, S.~J. Wright, \emph{Numerical
Optimization}, 2nd ed., Springer, 2006.
\bibitem{pff:golub-vanloan} G.~H. Golub, C.~F. Van Loan, \emph{Matrix
Computations}, 4th ed., Johns Hopkins University Press, 2013.
\bibitem{pff:ajjarapu-christy} V.~Ajjarapu, C.~Christy, The continuation power
flow: a tool for steady state voltage stability analysis, \emph{IEEE
Transactions on Power Systems}, vol. 7, no. 1, pp. 416--423, 1992.
\bibitem{pff:vancutsem-vournas} T.~Van Cutsem, C.~Vournas, \emph{Voltage
Stability of Electric Power Systems}, Kluwer Academic Publishers, 1998.
\bibitem{pff:kundur} P.~Kundur, \emph{Power System Stability and Control},
McGraw-Hill, 1994.
\bibitem{pff:hintermuller} M.~Hinterm\"uller, K.~Ito, K.~Kunisch, The
primal-dual active set strategy as a semismooth Newton method, \emph{SIAM
Journal on Optimization}, vol. 13, no. 3, pp. 865--888, 2002.
\bibitem{pff:tinney-walker} W.~F. Tinney, J.~W. Walker, Direct solutions of
sparse network equations by optimally ordered triangular factorization,
\emph{Proceedings of the IEEE}, vol. 55, no. 11, pp. 1801--1809, 1967.
\bibitem{pff:george-liu} A.~George, J.~W.~H. Liu, The evolution of the minimum
degree ordering algorithm, \emph{SIAM Review}, vol. 31, no. 1, pp. 1--19, 1989.
\bibitem{pff:amestoy-amd} P.~R. Amestoy, T.~A. Davis, I.~S. Duff, An approximate
minimum degree ordering algorithm, \emph{SIAM Journal on Matrix Analysis and
Applications}, vol. 17, no. 4, pp. 886--905, 1996.
\bibitem{pff:george-nested} A.~George, Nested dissection of a regular finite
element mesh, \emph{SIAM Journal on Numerical Analysis}, vol. 10, no. 2,
pp. 345--363, 1973.
\bibitem{pff:davis-suitesparse} T.~A. Davis, \emph{Direct Methods for Sparse
Linear Systems}, SIAM, 2006.
\bibitem{pff:iwamoto} S.~Iwamoto and Y.~Tamura, A load flow calculation method
for ill-conditioned power systems, \emph{IEEE Transactions on Power Apparatus
and Systems}, vol. PAS-100, no. 4, pp. 1736--1743, 1981.
\bibitem{pff:dacosta} V.~M. da Costa, N.~Martins, and J.~L.~R. Pereira,
Developments in the Newton Raphson power flow formulation based on current
injections, \emph{IEEE Transactions on Power Systems}, vol. 14, no. 4,
pp. 1320--1326, 1999.
\bibitem{pff:bolognani} S.~Bolognani and S.~Zampieri, On the existence and linear
approximation of the power flow solution in power distribution networks,
\emph{IEEE Transactions on Power Systems}, vol. 31, no. 1, pp. 163--172, 2016.
\bibitem{pff:simpsonporco} J.~W. Simpson-Porco, F.~D\"orfler, and F.~Bullo,
Voltage collapse in complex power grids, \emph{Nature Communications}, vol. 7,
art. 10790, 2016.
\end{thebibliography}}
